\documentclass[10pt,a4paper]{amsart}
\usepackage{amsmath,amssymb,amsthm,graphicx}
% For dealing with references we use the comment environment
\usepackage{verbatim}
\newenvironment{reference}{\comment}{\endcomment}
%\newenvironment{reference}{}{}
\newenvironment{slogan}{\comment}{\endcomment}
\newenvironment{history}{\comment}{\endcomment}
% For commutative diagrams we use Xy-pic
\usepackage[all]{xy}

% We use 2cell for 2-commutative diagrams.
\xyoption{2cell}
\UseAllTwocells
\theoremstyle{plain}
\newtheorem{theorem}[subsection]{Theorem}
\newtheorem{proposition}[subsection]{Proposition}
\newtheorem{lemma}[subsection]{Lemma}

\theoremstyle{definition}
\newtheorem{definition}[subsection]{Definition}
\newtheorem{example}[subsection]{Example}
\newtheorem{exercise}[subsection]{Exercise}
\newtheorem{situation}[subsection]{Situation}

\theoremstyle{remark}
\newtheorem{remark}[subsection]{Remark}
\newtheorem{remarks}[subsection]{Remarks}

\numberwithin{equation}{subsection}
% Macros
%
\def\lim{\mathop{\mathrm{lim}}\nolimits}
\def\colim{\mathop{\mathrm{colim}}\nolimits}
\def\Spec{\mathop{\mathrm{Spec}}}
\def\Hom{\mathop{\mathrm{Hom}}\nolimits}
\def\Ext{\mathop{\mathrm{Ext}}\nolimits}
\def\SheafHom{\mathop{\mathcal{H}\!\mathit{om}}\nolimits}
\def\SheafExt{\mathop{\mathcal{E}\!\mathit{xt}}\nolimits}
\def\Sch{\mathit{Sch}}
\def\Mor{\mathop{\mathrm{Mor}}\nolimits}
\def\Ob{\mathop{\mathrm{Ob}}\nolimits}
\def\Sh{\mathop{\mathit{Sh}}\nolimits}
\def\NL{\mathop{N\!L}\nolimits}
\def\CH{\mathop{\mathrm{CH}}\nolimits}
\def\proetale{{pro\text{-}\acute{e}tale}}
\def\etale{{\acute{e}tale}}
\def\QCoh{\mathit{QCoh}}
\def\Ker{\mathop{\mathrm{Ker}}}
\def\Im{\mathop{\mathrm{Im}}}
\def\Coker{\mathop{\mathrm{Coker}}}
\def\Coim{\mathop{\mathrm{Coim}}}

% Boxtimes
%
\DeclareMathSymbol{\boxtimes}{\mathbin}{AMSa}{"02}

%
% Macros for moduli stacks/spaces
%
\def\QCohstack{\mathcal{QC}\!\mathit{oh}}
\def\Cohstack{\mathcal{C}\!\mathit{oh}}
\def\Spacesstack{\mathcal{S}\!\mathit{paces}}
\def\Quotfunctor{\mathrm{Quot}}
\def\Hilbfunctor{\mathrm{Hilb}}
\def\Curvesstack{\mathcal{C}\!\mathit{urves}}
\def\Polarizedstack{\mathcal{P}\!\mathit{olarized}}
\def\Complexesstack{\mathcal{C}\!\mathit{omplexes}}
% \Pic is the operator that assigns to X its picard group, usage \Pic(X)
% \Picardstack_{X/B} denotes the Picard stack of X over B
% \Picardfunctor_{X/B} denotes the Picard functor of X over B
\def\Pic{\mathop{\mathrm{Pic}}\nolimits}
\def\Picardstack{\mathcal{P}\!\mathit{ic}}
\def\Picardfunctor{\mathrm{Pic}}
\def\Deformationcategory{\mathcal{D}\!\mathit{ef}}
\usepackage{cleveref}
\crefname{originalreference}{original reference}{original references}
\Crefname{originalreference}{Original reference}{Original references}
\makeatletter
\newlabel{situation-bootstrap}{{0EFU}{source}{}{sourcepointer-1}{}}
\newlabel{situation-bootstrap@cref}{{[originalreference][0][]0EFU}{[1][0][]source}{}{}{}}
\makeatother
\crefname{originallemma}{original lemma}{original lemmas}
\Crefname{originallemma}{Original lemma}{Original lemmas}
\makeatletter
\newlabel{lemma-bootstrap-inherited}{{0EFV}{source}{}{sourcepointer-2}{}}
\newlabel{lemma-bootstrap-inherited@cref}{{[originallemma][0][]0EFV}{[1][0][]source}{}{}{}}
\makeatother
\crefname{originallemma}{original lemma}{original lemmas}
\Crefname{originallemma}{Original lemma}{Original lemmas}
\makeatletter
\newlabel{lemma-algebraize-local-cohomology-bis}{{0EFR}{source}{}{sourcepointer-3}{}}
\newlabel{lemma-algebraize-local-cohomology-bis@cref}{{[originallemma][0][]0EFR}{[1][0][]source}{}{}{}}
\makeatother
\crefname{originallemma}{original lemma}{original lemmas}
\Crefname{originallemma}{Original lemma}{Original lemmas}
\makeatletter
\newlabel{lemma-ML-local}{{0DX1}{source}{}{sourcepointer-4}{}}
\newlabel{lemma-ML-local@cref}{{[originallemma][0][]0DX1}{[1][0][]source}{}{}{}}
\makeatother
\crefname{originallemma}{original lemma}{original lemmas}
\Crefname{originallemma}{Original lemma}{Original lemmas}
\makeatletter
\newlabel{local-cohomology-lemma-zero}{{0EF9}{source}{}{sourcepointer-5}{}}
\newlabel{local-cohomology-lemma-zero@cref}{{[originallemma][0][]0EF9}{[1][0][]source}{}{}{}}
\makeatother
\crefname{originallemma}{original lemma}{original lemmas}
\Crefname{originallemma}{Original lemma}{Original lemmas}
\makeatletter
\newlabel{local-cohomology-lemma-essential-image}{{0EFA}{source}{}{sourcepointer-6}{}}
\newlabel{local-cohomology-lemma-essential-image@cref}{{[originallemma][0][]0EFA}{[1][0][]source}{}{}{}}
\makeatother
\crefname{originallemma}{original lemma}{original lemmas}
\Crefname{originallemma}{Original lemma}{Original lemmas}
\makeatletter
\newlabel{dualizing-lemma-universally-catenary}{{0A80}{source}{}{sourcepointer-7}{}}
\newlabel{dualizing-lemma-universally-catenary@cref}{{[originallemma][0][]0A80}{[1][0][]source}{}{}{}}
\makeatother
\crefname{originalsection}{original section}{original sections}
\Crefname{originalsection}{Original section}{Original sections}
\makeatletter
\newlabel{dualizing-section-dimension-function}{{0A7W}{source}{}{sourcepointer-8}{}}
\newlabel{dualizing-section-dimension-function@cref}{{[originalsection][0][]0A7W}{[1][0][]source}{}{}{}}
\makeatother
\crefname{originallemma}{original lemma}{original lemmas}
\Crefname{originallemma}{Original lemma}{Original lemmas}
\makeatletter
\newlabel{dualizing-lemma-special-case-local-duality}{{0AAK}{source}{}{sourcepointer-9}{}}
\newlabel{dualizing-lemma-special-case-local-duality@cref}{{[originallemma][0][]0AAK}{[1][0][]source}{}{}{}}
\makeatother
\title{Lemma bootstrap in algebraization.tex}
\author{The Stacks Project Authors}
\date{stacks-project:a04446e57ec1fbc252a871afcec7752fb2807b14; GNU Free Documentation License 1.2 (as published in the Stacks Project COPYING)}
\begin{document}
\maketitle
\noindent\textit{Editorial note.} All mathematical statement and proof text is reproduced verbatim from the author's chapter algebraization.tex of the Stacks Project, fixed at commit a04446e57ec1fbc252a871afcec7752fb2807b14. The pinned preamble is supplied as its own file; the wrapper class, the xr-hyper wiring and the \textbackslash{}externaldocument list are omitted because the assembler supplies the class.\par\medskip
\noindent\textit{Editorial note.} Numbering inside this standalone excerpt is editorial, not the full Stacks book print numbering.\par\medskip
\noindent\textit{Editorial note.} The candidate is from the corrected band-I follow-on screen (line C) of the extractor's stacks-candidates line: its author label is absent from the live master in every form searched and it belongs to none of the earlier judged pools. The H2 judgement and its reason were made in this module by reading the author's own statement and complete proof in the fixed checkout. Original candidate key: stacks:algebraization.tex:lemma-bootstrap.\par\medskip
\noindent\textit{Editorial note.} Reference labels in the proof resolve to navigation pointers to the original author source; these pointers do not reproduce proofs of the prerequisites.\par\medskip
\noindent\textit{Allowed standard prerequisites (editorial).}\begin{itemize}
\item Prerequisites named by the author inside the counted span. The counted span is reproduced verbatim, including the author's own citations. Its prerequisites are the earlier results of the same fixed Stacks Project version that the author cites there and in the chapter's lead-in.
\end{itemize}\medskip
\ifdefined\hypertarget\hypertarget{sourcepointer-1}{}\fi
\noindent\textit{Source reference.} Original reference 0EFU: sources/algebraization.tex, source lines 2646--2646. Navigation pointer used in the counted span; the span contains the pointer only and the cited result is not reproduced inside it. Registered as a non-load-bearing navigation stub in the established form for Stacks material. No bibliography citation is involved.\par\medskip
\ifdefined\hypertarget\hypertarget{sourcepointer-2}{}\fi
\noindent\textit{Source reference.} Original lemma 0EFV: sources/algebraization.tex, source lines 2727--2727. Navigation pointer used in the counted span; the span contains the pointer only and the cited result is not reproduced inside it. Registered as a non-load-bearing navigation stub in the established form for Stacks material. No bibliography citation is involved.\par\medskip
\ifdefined\hypertarget\hypertarget{sourcepointer-3}{}\fi
\noindent\textit{Source reference.} Original lemma 0EFR: sources/algebraization.tex, source lines 2922--2922. Navigation pointer used in the counted span; the span contains the pointer only and the cited result is not reproduced inside it. Registered as a non-load-bearing navigation stub in the established form for Stacks material. No bibliography citation is involved.\par\medskip
\ifdefined\hypertarget\hypertarget{sourcepointer-4}{}\fi
\noindent\textit{Source reference.} Original lemma 0DX1: sources/algebraization.tex, source lines 664--664. Navigation pointer used in the counted span; the span contains the pointer only and the cited result is not reproduced inside it. Registered as a non-load-bearing navigation stub in the established form for Stacks material. No bibliography citation is involved.\par\medskip
\ifdefined\hypertarget\hypertarget{sourcepointer-5}{}\fi
\noindent\textit{Source reference.} Original lemma 0EF9: sources/local-cohomology.tex, source lines 1049--1049. Navigation pointer used in the counted span; the span contains the pointer only and the cited result is not reproduced inside it. Registered as a non-load-bearing navigation stub in the established form for Stacks material. No bibliography citation is involved.\par\medskip
\ifdefined\hypertarget\hypertarget{sourcepointer-6}{}\fi
\noindent\textit{Source reference.} Original lemma 0EFA: sources/local-cohomology.tex, source lines 1100--1100. Navigation pointer used in the counted span; the span contains the pointer only and the cited result is not reproduced inside it. Registered as a non-load-bearing navigation stub in the established form for Stacks material. No bibliography citation is involved.\par\medskip
\ifdefined\hypertarget\hypertarget{sourcepointer-7}{}\fi
\noindent\textit{Source reference.} Original lemma 0A80: sources/dualizing.tex, source lines 3324--3324. Navigation pointer used in the counted span; the span contains the pointer only and the cited result is not reproduced inside it. Registered as a non-load-bearing navigation stub in the established form for Stacks material. No bibliography citation is involved.\par\medskip
\ifdefined\hypertarget\hypertarget{sourcepointer-8}{}\fi
\noindent\textit{Source reference.} Original section 0A7W: sources/dualizing.tex, source lines 3224--3224. Navigation pointer used in the counted span; the span contains the pointer only and the cited result is not reproduced inside it. Registered as a non-load-bearing navigation stub in the established form for Stacks material. No bibliography citation is involved.\par\medskip
\ifdefined\hypertarget\hypertarget{sourcepointer-9}{}\fi
\noindent\textit{Source reference.} Original lemma 0AAK: sources/dualizing.tex, source lines 3510--3510. Navigation pointer used in the counted span; the span contains the pointer only and the cited result is not reproduced inside it. Registered as a non-load-bearing navigation stub in the established form for Stacks material. No bibliography citation is involved.\par\medskip
\begin{lemma}
\label{lemma-bootstrap}
In Situation \ref{situation-bootstrap} the inverse systems
$\{H^i_T(I^nM)\}_{n \geq 0}$ are pro-zero for $i \leq s$.
Moreover, there exists an integer $m_0$ such that for all
$m \geq m_0$ there exists an integer $m'(m) \geq m$ such that for
$k \geq m'(m)$ the image of
$H^{s + 1}_T(I^kM) \to H^{s + 1}_T(I^mM)$
maps injectively to $H^{s + 1}_T(I^{m_0}M)$.
\end{lemma}

\begin{proof}
Fix $m$. Let $\mathfrak q \in T$.
By Lemmas \ref{lemma-bootstrap-inherited} and
\ref{lemma-algebraize-local-cohomology-bis}
we see that
$$
H^i_\mathfrak q(M_\mathfrak q)
\longrightarrow
\lim H^i_\mathfrak q(M_\mathfrak q/I^nM_\mathfrak q)
$$
is an isomorphism for $i \leq s$. The inverse systems
$\{H^i_\mathfrak q(I^nM_\mathfrak q)\}_{n \geq 0}$ and
$\{H^i_\mathfrak q(M/I^nM)\}_{n \geq 0}$
satisfy the Mittag-Leffler condition for all $i$, see
Lemma \ref{lemma-ML-local}. Thus looking at the inverse system of
long exact sequences
$$
0 \to H^0_\mathfrak q(I^nM_\mathfrak q) \to
H^0_\mathfrak q(M_\mathfrak q) \to
H^0_\mathfrak q(M_\mathfrak q/I^nM_\mathfrak q) \to
H^1_\mathfrak q(I^nM_\mathfrak q) \to
H^1_\mathfrak q(M_\mathfrak q) \to \ldots
$$
we conclude (some details omitted) that there exists an integer
$m'(m, \mathfrak q) \geq m$ such that for all $k \geq m'(m, \mathfrak q)$
the map
$H^i_\mathfrak q(I^kM_\mathfrak q) \to H^i_\mathfrak q(I^mM_\mathfrak q)$
is zero for $i \leq s$ and the image of
$H^{s + 1}_\mathfrak q(I^kM_\mathfrak q) \to
H^{s + 1}_\mathfrak q(I^mM_\mathfrak q)$
is independent of $k \geq m'(m, \mathfrak q)$ and
maps injectively into $H^{s + 1}_\mathfrak q(M_\mathfrak q)$.

\medskip\noindent
Suppose we can show that $m'(m, \mathfrak q)$ can be chosen
independently of $\mathfrak q \in T$.
Then the lemma follows immediately from
Local Cohomology, Lemmas \ref{local-cohomology-lemma-zero} and
\ref{local-cohomology-lemma-essential-image}.

\medskip\noindent
Let $\omega_A^\bullet$ be a dualizing complex. Let
$\delta : \Spec(A) \to \mathbf{Z}$ be the corresponding
dimension function. Recall that $\delta$ attains only a
finite number of values, see
Dualizing Complexes, Lemma \ref{dualizing-lemma-universally-catenary}.
Claim: for each $d \in \mathbf{Z}$ the integer
$m'(m, \mathfrak q)$ can be chosen independently
of $\mathfrak q \in T$ with $\delta(\mathfrak q) = d$.
Clearly the claim implies the lemma by what we said above.

\medskip\noindent
Pick $\mathfrak q \in T$ with $\delta(\mathfrak q) = d$.
Consider the ext modules
$$
E(n, j) = \text{Ext}^j_A(I^nM, \omega_A^\bullet)
$$
A key feature we will use is that these are finite $A$-modules.
Recall that $(\omega_A^\bullet)_\mathfrak q[-d]$ is a normalized
dualizing complex for $A_\mathfrak q$ by definition of the
dimension function associated to a dualizing complex, see
Dualizing Complexes, Section \ref{dualizing-section-dimension-function}.
The local duality theorem (Dualizing Complexes, Lemma
\ref{dualizing-lemma-special-case-local-duality}) tells us that
the $\mathfrak qA_\mathfrak q$-adic completion of
$E(n, -d - i)_\mathfrak q$ is Matlis dual to
$H^i_\mathfrak q(I^nM_\mathfrak q)$. Thus the choice of
$m'(m, \mathfrak q)$ for $i \leq s$ in the first paragraph tells us that
for $k \geq m'(m, \mathfrak q)$ and $j \geq -d - s$ the map
$$
E(m, j)_\mathfrak q \to E(k, j)_\mathfrak q
$$
is zero. Since these modules are finite and nonzero only
for a finite number of possible $j$ (small detail omitted),
we can find an open neighbourhood $W \subset \Spec(A)$ of $\mathfrak q$
such that
$$
E(m, j)_{\mathfrak q'} \to E(m'(m, \mathfrak q), j)_{\mathfrak q'}
$$
is zero for $j \geq -d - s$ for all $\mathfrak q' \in W$.
Then of course the maps $E(m, j)_{\mathfrak q'} \to E(k, j)_{\mathfrak q'}$
for $k \geq m'(m, \mathfrak q)$ are zero as well.

\medskip\noindent
For $i = s + 1$ corresponding to $j = - d - s - 1$ we obtain
from local duality and the results of the first paragraph that
$$
K_{k, \mathfrak q} =
\Ker(E(m, -d - s - 1)_\mathfrak q \to E(k, -d - s - 1)_\mathfrak q)
$$
is independent of $k \geq m'(m, \mathfrak q)$ and that
$$
E(0, -d - s - 1)_\mathfrak q \to
E(m, -d - s - 1)_\mathfrak q/K_{m'(m, \mathfrak q), \mathfrak q}
$$
is surjective. For $k \geq m'(m, \mathfrak q)$ set
$$
K_k = \Ker(E(m, -d - s - 1) \to E(k, -d - s - 1))
$$
Since $K_k$ is an increasing sequence of submodules of the finite
module $E(m, -d - s - 1)$ we see that, at the cost of increasing
$m'(m, \mathfrak q)$ a little bit, we may assume
$K_{m'(m, \mathfrak q)} = K_k$ for $k \geq m'(m, \mathfrak q)$.
After shrinking $W$ further if necessary, we may also assume that
$$
E(0, -d - s - 1)_{\mathfrak q'} \to
E(m, -d - s - 1)_{\mathfrak q'}/K_{m'(m, \mathfrak q), \mathfrak q'}
$$
is surjective for all $\mathfrak q' \in W$ (as before use that
these modules are finite
and that the map is surjective after localization at $\mathfrak q$).

\medskip\noindent
Any subset, in particular
$T_d = \{\mathfrak q \in T \text{ with }\delta(\mathfrak q) = d\}$,
of the Noetherian topological space $\Spec(A)$
with the endowed topology is Noetherian and hence quasi-compact.
Above we have seen that for every $\mathfrak q \in T_d$
there is an open neighbourhood $W$ where
$m'(m, \mathfrak q)$ works for all $\mathfrak q' \in T_d \cap W$.
We conclude that we can find an integer $m'(m, d)$ such that for all
$\mathfrak q \in T_d$ we have
$$
E(m, j)_\mathfrak q \to E(m'(m, d), j)_\mathfrak q
$$
is zero for $j \geq -d - s$ and with
$K_{m'(m, d)} = \Ker(E(m, -d - s - 1) \to E(m'(m, d), -d - s - 1))$
we have
$$
K_{m'(m, d), \mathfrak q} =
\Ker(E(m, -d - s - 1)_{\mathfrak q} \to E(k, -d - s - 1)_{\mathfrak q})
$$
for all $k \geq m'(m, d)$ and the map
$$
E(0, -d - s - 1)_\mathfrak q \to
E(m, -d - s - 1)_\mathfrak q/K_{m'(m, d), \mathfrak q}
$$
is surjective. Using the local duality theorem again (in the opposite
direction) we conclude that the claim is correct. This finishes the proof.
\end{proof}

\end{document}
